% Hardware, software, and the setup of the Day 12 lab.
% The values come from Labs/day12/README.md.
\begin{frame}{Hardware and software}
  \begin{columns}[T]
    \begin{column}{0.47\textwidth}
      \small
      \textbf{Hardware for each group}
      \begin{itemize}
        \item XIAOML Kit and a USB-C cable.
        \item Raspberry Pi 5 with the active cooler, the card of your group,
              and the 27 W power supply.
        \item Laptop in the lab network \code{edgeai-lab}.
        \item For the class: the instructor laptop with a Mosquitto broker
              at \code{X.10}.
      \end{itemize}
    \end{column}
    \begin{column}{0.49\textwidth}
      \small
      \textbf{Software}
      \begin{itemize}
        \item Card: Mosquitto and its clients, the environment
              \code{\textasciitilde/yolo} with \code{paho-mqtt}.
        \item Laptop: Python 3 with \code{paho-mqtt} 2.x, the Arduino IDE
              with the esp32 core 3.3.12, PubSubClient 2.8, U8g2.
        \item The ZIP library of your Day 5 keyword project.
        \item MicroPython v1.29.0, \code{esptool}, \code{mpremote}
              (HW-01).
      \end{itemize}
    \end{column}
  \end{columns}
  \vskip 0.4em
  \small
  \renewcommand{\arraystretch}{1.15}
  \begin{tabular}{@{}L{0.24\textwidth}L{0.70\textwidth}@{}}
    Lab files & \code{Labs/day12/} on the laptop, and a copy in
                \code{\textasciitilde/edgeai/} on the Raspberry Pi. \\
    Names     & \code{gNN} is your group: \code{pi-NN},
                \code{X.(100+NN)}. The examples use \code{g07} and \code{X}
                = \code{192.168.8}. \\
  \end{tabular}
\end{frame}

\begin{frame}{The setup}
  \begingroup\centering
  \begin{tikzpicture}[font=\scriptsize, >=stealth,
      box/.style={draw=coolgray, rounded corners=2pt, align=center,
                  minimum height=0.8cm, text width=1.8cm}]
    \draw[darkcyan, dashed, rounded corners=3pt] (1.45,-1.35) rectangle
      (6.25,1.15);
    \node[darkcyan, anchor=south west] at (1.5,1.15) {Raspberry Pi
      (\texttt{pi-07})};
    \node[box, fill=cardinalred!12, text width=1.4cm] (x) at (-0.45,0.5)
      {\textbf{XIAOML Kit}\\ keyword model};
    \node[box, fill=boxgray] (lb) at (2.65,0.5) {Mosquitto};
    \node[box, fill=boxgray] (d) at (2.65,-0.75) {\texttt{decide.py}};
    \node[box, fill=boxgray] (f) at (5.05,0.5) {\texttt{forwarder.py}};
    \node[box, fill=boxgray] (q) at (5.05,-0.75) {queue on disk};
    \node[box, fill=darkcyan!20, text width=1.7cm] (c) at (8.3,0.5)
      {\textbf{Instructor laptop}\\ cloud broker};
    \draw[->] (x.east |- 0,0.65) -- node[above] {events} (lb.west |- 0,0.65);
    \draw[->, darkcyan] (lb.west |- 0,0.35) -- node[below] {\texttt{cmd}}
      (x.east |- 0,0.35);
    \draw[<->] (lb) -- (d);
    \draw[->] (lb) -- (f);
    \draw[<->] (f) -- (q);
    \draw[->] (f) -- node[above] {link} (c);
    \node[cardinalred, font=\normalsize] at (6.8,0.48) {$\times$};
  \end{tikzpicture}\par\endgroup
  \vskip 0.4em
  \small
  \begin{itemize}
    \item The XIAO, the broker, and \code{decide.py} form the local system.
          The LED command needs no link to the laptop.
    \item \code{forwarder.py} writes each decision and each summary to a
          queue on the disk, and sends it to the cloud broker with QoS 1.
    \item On the group laptop, \code{cloud\_check.py} counts the messages of
          the group on the cloud broker.
  \end{itemize}
\end{frame}

\begin{frame}{The topics of the lab}
  \footnotesize
  \renewcommand{\arraystretch}{1.12}
  \setlength{\tabcolsep}{4pt}
  \begin{tabular}{@{}L{3.6cm}L{1.6cm}L{3.2cm}L{1.3cm}@{}}
    \toprule
    \textbf{Topic} & \textbf{From} & \textbf{Payload} & \textbf{QoS,
      retained} \\
    \midrule
    \texttt{edgeai/g07/xiao/imu} & XIAO (B1) & one IMU reading each 200 ms
      & 0, no \\
    \texttt{edgeai/g07/xiao/result} & XIAO & one keyword event & 0, no \\
    \texttt{edgeai/g07/xiao/stats} & XIAO & each 10 s: windows, events,
      memory & 0, no \\
    \texttt{edgeai/g07/xiao/status} & XIAO & \texttt{online}, will
      \texttt{offline} & 1, yes \\
    \texttt{edgeai/g07/xiao/cmd} & Raspberry Pi & \texttt{led=1},
      \texttt{led=0} & 1, no \\
    \texttt{edgeai/g07/pi/decision} & Raspberry Pi & one change of state &
      1, no \\
    \texttt{edgeai/g07/pi/summary} & Raspberry Pi & each 10 s: the state and
      the counts & 1, no \\
    \texttt{edgeai/g07/pi/status} & Raspberry Pi & \texttt{online}, will
      \texttt{offline} & 1, yes \\
    \bottomrule
  \end{tabular}
  \vskip 0.4em
  \small
  The forwarder sends \code{pi/decision}, \code{pi/summary}, and
  \code{xiao/stats} to the cloud. Each message with JSON has a sequence
  number \code{seq}.
\end{frame}
